\section{Declaración del alcance}
Este proyecto se centró en diseñar, desarrollar y poner en marcha los módulos Aprender a Aprender y Escucha Activa para el prototipo de plataforma gamificada orientada a la enseñanza de habilidades blandas. El trabajo estuvo dirigido específicamente a estudiantes de Ingeniería de Sistemas de la Universidad del Valle, sede Tuluá, con el propósito de fortalecer su capacidad de aprendizaje autónomo y promover prácticas de escucha activa mediante experiencias interactivas. Más allá del componente técnico, se buscó enfatizar la necesidad de integrar habilidades transversales dentro de la formación académica y profesional de ingeniería.
\\
\\
El desarrollo del proyecto enfrentó limitaciones asociadas tanto al tiempo como a la disponibilidad de recursos humanos, dado que su construcción fue realizada por una sola persona dentro de un periodo máximo de 12 meses. En consecuencia, el alcance se concentró en dos habilidades blandas específicaas (aprender a aprender y escucha activa) dejando fuera otros módulos potenciales. 
\\
\\
Esta investigación se basó en condiciones necesarias para garantizar su ejecución: el interés y la participación activa de los estudiantes en el uso de la plataforma, el acompañamiento docente durante todas las etapas del desarrollo, y el acceso a recursos institucionales adecuados para la revisión, prueba y seguimiento del proyecto. Estas condiciones humanas y organizacionales permitieron avanzar en el proceso de implementación y sentaron las bases para la validación del prototipo.
\\
\\
A partir de estas condiciones iniciales, la ejecución del proyecto también dependió de requisitos técnicos fundamentales para asegurar su éxito: el acceso continuo a servicios externos como modelos de procesamiento de lenguaje natural (PLN) y sintetizadores de voz; la disponibilidad de internet para probar la plataforma de manera constante; un entorno de desarrollo con los recursos necesarios para ejecutar e integrar los servicios basados en LLM; y la utilización de herramientas modernas de desarrollo web para construir tanto el frontend como el backend del prototipo. Estas condiciones resultaron esenciales para la implementación y el correcto funcionamiento de los módulos desarrollados, especialmente aquellos que requieren generación dinámica de contenido.
\\
\\
El alcance del proyecto incluyó los siguientes componentes principales:
\begin{enumerate}
    \item \textbf{Módulo de gestión del aprendizaje (Aprender a Aprender):}
    Este módulo integró herramientas para la planificación, organización y seguimiento del aprendizaje mediante la creación de metas, objetivos y tareas. Incluyó un tablero Kanban para la gestión visual de actividades y un sistema de autoevaluaciones que promueva la reflexión del estudiantes. 

    \item \textbf{Generación y edición de roadmaps de aprendizaje:}
    Se desarrolló un editor visual que permite a los usuarios construir, organizar y ajustar rutas de aprendizajes personalizadas. Este componente incorporó un agente conversacional basado en modelos de procesamiento de lenguaje natural (PLN), capaz de generar roadmaps adaptados a las necesidades e intereses del estudiante, facilitando la orientación inicial del proceso de aprendizaje.

    \item \textbf{Integración bidireccional entre roadmaps y planificación del aprendizaje:}
    La plataforma permitió transformar la estructura de metas, objetivos y tareas en un roadmap visual que podía hacerse público para compartir distintas formas de abordar un proceso de aprendizaje. A su vez, los roadmpas existentes podían convertirse nuevamente en metas, objetivos y tareas, de modo que cada usuario pudiera adoptar o adaptar rutas creadas por otros según sus necesidades.

    \item \textbf{Módulo de entrenamiento en escucha activa (gamificado):}
    Este módulo ofreció una experiencia interactiva compuesta por sesiones de juego divididas en rondas dinámicas. Cada ronda presentó desafíos auditivos generados mediante inteligencia artificial, con diferentes modos de juego (selección múltiple, completar espacios, parafraseo, resumen y formulación de preguntas de clarificación). El sistema evaluó automáticamente las respuestas y ofreció retroalimentación inmediata para fortalecer la comprensión auditiva y la atención activa.

    \item \textbf{Generación y reproducción de audio con modelos de IA:}
    Los desafíos del módulo de escucha activa integraron síntesis de voz mediante servicios externos, lo que permitió obtener audios con distintos estilos, voces y niveles de dificultad. El sistema gestionó la creación, almacenamiento y reproducción de estos recursos para garantizar la disponibilidad del contenido auditivo requerido en cada ronda de juego.
\end{enumerate}